\begin{textsample}{POS Dim 6 – mg – Score 7.00 – mg/mg\_em\_108.txt}  \label{ex:f6_pos_246}
5.6 A Educação Profissional e Técnica e os Eixos \textbf{Estruturantes} Em um mundo cada vez mais interligado, conectado e estruturado pela \textbf{tecnologia} \textbf{digital}, são exigidas dos jovens novas habilidades, de modo que possam se adaptar com maior rapidez às mudanças ocorridas na sociedade, respondendo, de modo assertivo, aos desafios apresentados nas diversas esferas de atuação.
A transformação nos processos de produção, nas formas de construção do conhecimento, nos meios de comunicação e, principalmente, as exigências de novas competências e habilidades para atuar no mundo de forma proficiente mostram o quanto a formação escolar precisa se reorganizar a partir de princípios inovadores e flexíveis.
Em um estudo recente, denominado “ O futuro do trabalho ” ( 2016 ), o Fórum Econômico Mundial publicou quais habilidades serão as mais importantes no mercado de trabalho no ano de 2020, estabelecendo um comparativo com as que eram consideradas significativas em 2015.
Podemos perceber que tem sido necessário menos de uma década para que as demandas \textbf{formativas} sofram impactos derivados das transformações ocorridas na sociedade e nos meios de consumo e produção.
Os documentos curriculares precisam, portanto, estar alinhados a esse ritmo, sendo permeáveis às mudanças e às novas demandas sociais.
Em alinhamento ao apresentado, o Currículo Referência será composto por uma Formação Geral Básica em que os estudantes desenvolverão, fortalecerão e ampliarão os conhecimentos referentes às áreas do conhecimento e por Itinerários \textbf{Formativos}, cuja oferta se estruturará a partir da capacidade e das possibilidades do sistema de ensino, das características e singularidades do território e do contexto local em que a escola está inserida.
Considerará, também, no caso da oferta do 5º itinerário, dados referentes ao mundo do trabalho, tais como indicadores de demanda dos setores produtivos regionais e empregabilidade.
A implantação da Educação Profissional e Tecnológica como itinerário \textbf{formativo} será efetivada a partir de algumas opções de composição dos itinerários, contemplando sempre tempos e espaços para a construção e o fortalecimento do Projeto de vida e o desenvolvimento de competências e habilidades relacionadas à preparação geral para o mundo do trabalho.
Os percursos poderão contemplar também a oferta de cursos técnicos, de acordo com o CNCT, de cursos FIC organizados de forma articulada e até mesmo a partir da estruturação para a oferta de um Programa de aprendizagem pela Rede estadual de educação de Minas Gerais.
O percurso \textbf{formativo} da preparação básica para o trabalho se organiza a partir da abordagem e da integração dos diferentes eixos \textbf{estruturantes}, sendo que as habilidades a eles associadas somam -se a outras habilidades básicas requeridas indistintamente pelo mundo do trabalho e as habilidades específicas requeridas pelas distintas ocupações.
Ou seja, ao escolher o itinerário \textbf{formativo} de Formação Técnica e Profissional, necessariamente, os estudantes devem vivenciar um módulo de Preparação Básica para o Trabalho que se apresenta a partir dos seguintes eixos \textbf{estruturantes} : A investigação científica que, aplicada à formação técnica e profissional, proporciona o embasamento teórico e fenomenológico necessário para estruturar os conceitos práticos e aplicáveis em todos os aspectos técnicos.
Sem o alicerce que a investigação científica proporciona, o ensino puramente técnico fica comprometido, resultando em profissionais que apenas reproduzem ensinos metódicos, sem se preocupar com a origem do processo e a importância do que está sendo produzido.
Os processos criativos que possibilitam ao estudante desenvolver um pensamento abstrato com o propósito de resolver \textbf{problemas} por meio de modelos, experimentos e protótipos que apresentam soluções criativas, simples e de baixo custo.
Para o desenvolvimento dos processos criativos é necessária a utilização de instrumentos modernos, como Design Thinking ( Metodologia usada para solução de \textbf{problemas} ), como eixo norteador da construção de uma visão que busque se apoiar em estrutura de processos que auxiliam em escolhas mais assertivas para resolução de \textbf{problemas}.
Em um mundo tão interligado, em que as informações são adquiridas quase instantaneamente, torna -se viável que o ensino técnico também seja reestruturado a partir de uma visão empática e humana.
Portanto, não basta apenas saber e conhecer técnicas, mas aplicar o conhecimento visando ao bem-estar da sociedade em que o jovem está inserido.
A \textbf{mediação} e intervenção sociocultural, trazendo a preocupação em desenvolver soluções para gerir conflitos em seus vários âmbitos possíveis e proporcionar o desenvolvimento de inteligência emocional, buscando propiciar a construção de ambientes saudáveis e harmônicos.
O \textbf{empreendedorismo}, visando à construção de estratégias que, muito além de ensinar ao estudante práticas comerciais e planilhas eletrônicas para gestão financeira ou criação de networks, se pautem em uma formação técnica em que o jovem seja protagonista na tomada de suas decisões e entenda a necessidade de mudar suas expectativas de futuro.
O que cabe à escola é contribuir, por meio de atividades e cursos pertinentes, para o desenvolvimento de visão ampla de mercado e negócios, dar base para o aprimoramento e experiências relevantes no processo de aprendizagem do estudante, como também buscar parcerias que viabilizem a amplificação das oportunidades.

% matched lemmas: digital, empreendedorismo, estruturante, formativo, mediação, problema, tecnologia
\end{textsample}
